% Generated draft: requires mathematical and source review.
\documentclass[11pt]{article}
\usepackage[margin=25mm]{geometry}
\usepackage{fontspec}
\setmainfont{Latin Modern Roman}
\usepackage{amsmath,amssymb,mathtools,unicode-math}
\setmathfont{Latin Modern Math}
\usepackage{enumitem,needspace,xurl,hyperref}
\setlength{\parindent}{0pt}
\setlength{\parskip}{5pt}
\setlength{\emergencystretch}{3em}
\newcommand{\N}{\mathbb{N}}
\newcommand{\Z}{\mathbb{Z}}
\newcommand{\Q}{\mathbb{Q}}
\newcommand{\R}{\mathbb{R}}
\newcommand{\C}{\mathbb{C}}
\newcommand{\F}{\mathbb{F}}
\newcommand{\A}{\mathbb{A}}
\newcommand{\PP}{\mathbb{P}}
\newcommand{\E}{\mathbb{E}}
\newcommand{\eps}{\varepsilon}
\newcommand{\dd}{\,\mathrm{d}}
\begin{document}
\section*{UnsolvedMath ID: 30006541}
\textbf{Source problem code:} catalogue entry
\subsection{Definitions and mathematical statement}
The supplied record is a truncated passage beginning with the isolated word ``conjecture.'' It then states Theorems 3 and 4, rather than a standalone open problem. The following are transcriptions of those reported assertions, with their unresolved conventions made explicit.

For a finite nonempty simple undirected graph \(G\), the usual average degree is \(\overline d(G)=2|E(G)|/|V(G)|\). For \(S\subseteq V(G)\), the induced subgraph \(G[S]\) contains exactly the edges of \(G\) with both endpoints in \(S\). A subdivision of \(H\) replaces its edges by paths of positive length whose internal vertices are mutually disjoint and disjoint from the original vertices. An induced subdivision of \(H\) in \(G\) is an induced subgraph of \(G\) isomorphic to such a subdivision. Write \(C_4\) for the cycle on four vertices. Under the usual convention, \(C_4\)-free means containing no subgraph isomorphic to \(C_4\), whether induced or otherwise; bipartite means admitting a partition into two independent vertex sets. These definitions explain standard terminology; the fragment does not explicitly specify its graph conventions or define \((c,t)\)-sparsity.

\textbf{Reported Theorem 3.} For every \(c>0\) and every graph \(H\), there exists a constant \(C\) such that every graph \(G\) satisfying
\[
G\text{ is }(c,t)\text{-sparse},\qquad \overline d(G)\ge Ct
\]
contains an induced subdivision of \(H\).

\textbf{Reported Theorem 4.} For \(k\ge1\) and \(c>0\), there exists \(C=C(k,c)\) such that every graph \(G\) satisfying
\[
G\text{ is }(c,t)\text{-sparse},\qquad \overline d(G)\ge Ct
\]
has a vertex set \(S\subseteq V(G)\) for which \(G[S]\) is bipartite and \(C_4\)-free and
\[
\overline d\bigl(G[S]\bigr)\ge k.
\]

The domain and quantification of \(t\), the domain of \(k\) beyond \(k\ge1\), and the definition of \((c,t)\)-sparsity are absent. Accordingly, these are qualified transcriptions, not complete closed mathematical formulations. The passage describes Theorem 4 as a strengthening of an unavailable Theorem 2 and says Theorem 3 follows using a result attributed to K\"uhn and Osthus [19]: \(C_4\)-free graphs of sufficiently large average degree contain an induced subdivision of \(H\). The hypotheses and threshold dependence of that cited result are not supplied.

The registry labels the record ``open,'' but its supplied statement-verification field says the passage consists of two theorems proving a formerly conjectured statement and contains no remaining open problem. No separate conjecture or open target can therefore be recovered from the supplied evidence.
\subsection{Short English statement}
This fragment reports two theorems: sufficiently large average degree relative to \(t\) in a \((c,t)\)-sparse graph forces an induced subdivision of a prescribed graph, and also forces a bipartite \(C_4\)-free induced subgraph of prescribed average degree. Sparsity and parameter conventions are missing. The supplied verification says no open problem remains, despite the registry's ``open'' label; this draft preserves the reported assertions for review.
\subsection{Sources}
\begin{itemize}
\item Local UnsolvedMath registry, supplied complete record, permanent ID \(30006541\), exact source problem code \texttt{catalogue entry}, title ``Solved Sparse-Graph Conjecture Fragment.'' Collection: Oberwolfach Reports Open Problems; proposed year: 2026; no proposer is supplied. The record was updated at \texttt{2026-09-16T12:06:53Z}.
\item Supplied report citation and literature-source link: \emph{Combinatorics} (2026), Oberwolfach Reports, DOI 10.4171/owr/2026/1, pp. 5--79. \url{https://doi.org/10.4171/owr/2026/1}. This report is supplied as the authoritative provenance; the exact location within those pages and theorem authorship are not supplied.
\item The fragment attributes an auxiliary induced-subdivision result to K\"uhn and Osthus [19]. The bibliography entry for [19] is not supplied, so its publication details and full hypotheses cannot be reconstructed.
\item Literature assessment date: \texttt{2026-08-22}. This is supplied dated curation, not independently verified. It labels the record open while describing the extracted assertion as unresolved or unclear because the passage is truncated and contains no complete standalone problem claim. The background requests independent verification of time-sensitive status and cited claims.
\item The supplied \texttt{statement\_status} is \texttt{unrecoverable}. The supplied \texttt{statement\_verification} says the official report confirms a stray word attached to two theorems proving the formerly conjectured sparse-graph statement, with no remaining open problem; no clean statement was promoted, and the original record was retained publicly for provenance. This is a supplied verification claim, not a claim of independent source consultation here.
\end{itemize}
\end{document}
